\documentclass{../template/ffhsthesis}
\usepackage[a4paper, total={19cm, 26cm}]{geometry}
\usepackage[utf8]{inputenc} % this is needed for umlauts
\usepackage[ngerman]{babel} % this is needed for umlauts
\usepackage[T1]{fontenc}    % this is needed for correct output of umlauts in pdf
\usepackage{pgf,pgffor}
\usepackage{pgfplots}
\usepackage{float}%
\pgfplotsset{compat=newest}
\usepackage{calc}
\usepackage{verbatim}
\usepackage{csvsimple}
\usepackage{pdflscape}
\usepackage{multirow}
\usepackage{xfrac}


\makeatother

\setlength\parindent{0pt}




\begin{document}
% addional configuration from sam
\shorthandoff{"}


\dokumentTyp{Anhang der Bachelor-Thesis}
\studiengang{INF}
\title{Energie sparen durch bessere Software}
\subtitle{Analyse des Energieverbrauchs pro Befehlssatz} % optional
\titelbild[height=4.55cm,width=15cm]{images/Italy_Alps_and_Mediterranean.jpg}  % optional
\author{Samuel Riolo}
% \date{}
\wohnort{Kerzers}
%\referent{Name des Referenten\\ Titel\\Unterrichtetes Fach}
%\referent{Jürg Hofer}
\referent{Jürg Hofer\\ Dipl. El. Ing. ETH\\WING}
%\eingereichtBei{Prof.\ Dr.\ Martin Sutter\\Departement Informatik\\Departementsleiter} 


\maketitle



\startThesis
\chapter{Benchmark results with Raspberry Pi}
Das Raspberry Pi wurde in den folgenden Benchmarks mit einer konstante Spannung von \textbf{5.16 V} gespeist. Der zu testende Befehlssatz wurde \textbf{21,4 Milliarden} oder genau 21'474'836'480 mal ausgeführt.

\tikzset{every picture/.style=thick}





\section{Raw results}
\foreach \fileno in {0,...,100}{
\IfFileExists{raspberrydata/data\fileno.csv}{%
\InputIfFileExists{raspberrydata/data\fileno .tex}{}{}%



\section{\benchmark}
\description
\par
\textit{data\fileno .csv}

\begin{figure}[H]

\def\discon{$\wr\!\!\!\wr$}
\foreach \n  in {0,...,2}{%
\begin{minipage}{.33\textwidth}
%
\begin{tikzpicture}

    \begin{axis}[
            axis x line=middle,
            axis y line=middle,
            enlarge y limits=true,
            xmin=0, xmax=150000,
            ymin=425, ymax=445,
            width=6cm, height=4cm,     % size of the image
            grid = major,
            grid style={dashed, gray!30},
            ylabel=ampere (mA),
            xlabel=seconds (mS),
            extra x ticks = {1},
            xtick  = {50000, 100000},
            scaled x ticks = false, 
            x tick label style={/pgf/number format/.cd, fixed, fixed zerofill, int detect,1000 sep={\;},precision=3},  % Removing scientific notation
            extra y ticks = {425},
            ytick  = {430, 435, ..., 445},
            extra y tick style = {grid=none, yticklabel style={xshift=2.6mm, rotate=0, inner sep=.5pt}},
            extra y tick labels = {\includegraphics{almost_equals.eps}}]
            legend style={at={(0.1,-0.1)}, anchor=north}
         ]
          \addplot table [x=time\n, y=amp\n, col sep=comma, mark=none, ] {raspberrydata/data\fileno .csv};


    \end{axis}


\end{tikzpicture}%



% Example of an array in latex
%\def\names{{"5454.5","Frank","Laura","Joe"}}%
%\pgfmathparse{\names[0]}\pgfmathresult


\begin{tabular}{lcc}
Average:& \pgfmathparse{\average[\n]}\pgfmathresult mA\\
Median:& \pgfmathparse{\median[\n]}\pgfmathresult mA\\
Exec time: & \pgfmathparse{\exectime[\n]}\pgfmathresult mS\\
Exec order:& \pgfmathparse{\run[\n]}\pgfmathresult \\
\end{tabular}

\end{minipage}
}



\end{figure}


}
}


\begin{landscape}
\section{Results in tabular form}
\centering
\thispagestyle{empty}
\csvautotabular{raspberrytable.csv}
\par
\vspace{10pt}
\textbf{BM} = Benchmark, \textbf{T} = executiontime of an Bechmark, \textbf{Med} = Median (\sfrac{1}{2}-quantile)
\par
* The Energy (mW/) used for the benchmark with $\approx$21,4 billion cycles.
\end{landscape}



















\chapter{Benchmark results with Intel Galileo Gen2}




Das Galileo Gen2 Board wurde in den folgenden Benchmarks mit einer konstante Spannung von \textbf{6.36 V} gespeist. Der zu testende Befehlssatz wurde \textbf{6,4 Milliarden} oder genau 6'442'450'944 mal ausgeführt.

\tikzset{every picture/.style=thick}

\section{Raw results}

\foreach \fileno in {0,...,100}{
\IfFileExists{galileodata/data\fileno.csv}{%
\InputIfFileExists{galileodata/data\fileno .tex}{}{}%



\section{\benchmark}
\description
\par
\textit{data\fileno .csv}

\begin{figure}[H]

\def\discon{$\wr\!\!\!\wr$}
\foreach \n  in {0,...,2}{%
\begin{minipage}{.33\textwidth}
%
\begin{tikzpicture}

    \begin{axis}[
            axis x line=middle,
            axis y line=middle,
            enlarge y limits=true,
            xmin=0, xmax=150000,
            ymin=415, ymax=440,
            width=6cm, height=4cm,     % size of the image
            grid = major,
            grid style={dashed, gray!30},
            ylabel=ampere (mA),
            xlabel=seconds (mS),
            extra x ticks = {1},
            xtick  = {50000, 100000},
            scaled x ticks = false, 
            x tick label style={/pgf/number format/.cd, fixed, fixed zerofill, int detect,1000 sep={\;},precision=3},  % Removing scientific notation
            extra y ticks = {415},
            ytick  = {420, 425, ..., 440},
            extra y tick style = {grid=none, yticklabel style={xshift=2.6mm, rotate=0, inner sep=.5pt}},
            extra y tick labels = {\includegraphics{almost_equals.eps}}]
            legend style={at={(0.1,-0.1)}, anchor=north}
         ]
          \addplot table [x=time\n, y=amp\n, col sep=comma, mark=none, ] {galileodata/data\fileno .csv};


    \end{axis}


\end{tikzpicture}%



% Example of an array in latex
%\def\names{{"5454.5","Frank","Laura","Joe"}}%
%\pgfmathparse{\names[0]}\pgfmathresult


\begin{tabular}{lcc}
Average:& \pgfmathparse{\average[\n]}\pgfmathresult mA\\
Median:& \pgfmathparse{\median[\n]}\pgfmathresult mA\\
Exec time: & \pgfmathparse{\exectime[\n]}\pgfmathresult mS\\
Exec order:& \pgfmathparse{\run[\n]}\pgfmathresult \\
\end{tabular}

\end{minipage}
}



\end{figure}


}
}






\begin{landscape}
\section{Results in tabular form}
\centering
\thispagestyle{empty}
\csvautotabular{galileotable.csv}
\par
\vspace{10pt}
\textbf{BM} = Benchmark, \textbf{T} = executiontime of an Bechmark, \textbf{Med} = Median (\sfrac{1}{2}-quantile)
\par
* The Energy (mW/) used for the benchmark with $\approx$6.4 billion cycles.
\end{landscape}





\end{document}
